\section{Carrier Refinement}
\label{sec:carrier}

Beam Fusion is deliberately \emph{not} a drop-in initializer.  Its means are accurate early,
but covariance, opacity, and appearance start from geometry-only reasoning; dropped directly
into a standard 3DGS optimizer with densification, the optimizer tends to \emph{replace}
the initialization rather than refine it --- and the experiment then measures densification's
ability to recover, not the initializer.  Beam outputs are therefore treated as
\emph{carriers}: they pass through a short, fully compact-only maturation before any
comparison.  The accepted carrier sequence was selected by three preregistered, paired-root,
independently audited ablations\evidence{docs/ADR-002-carrier-refinement.md;
ara/logic/claims.md C30, C31} and consists of five stages:
\begin{enumerate}[leftmargin=1.6em,itemsep=1pt]
  \item fit-only Beam Fusion from reconstruction inputs ($K \ge 3$, no free primitive
        birth);
  \item renderer-aware covariance repair (\cref{sec:cov-repair});
  \item 380 compact fixed-topology SH0 updates with means, rotations, scales, opacity, and
        SH0 all trainable;
  \item strict projected-center pruning against every fitting view (\cref{sec:containment});
  \item 380 further compact SH0 updates with means frozen, so the pruning invariant cannot
        be undone.
\end{enumerate}
No clone, split, insertion, densification, opacity reset, higher-SH phase, or dense handover
exists on this path, and its entry point accepts no image-bearing argument.

\subsection{Renderer-aware covariance repair}
\label{sec:cov-repair}

The point renderer projects a 3D covariance as $P_{iv} = J_{iv} \Sigma_i J_{iv}\T + 0.3\,
I_2$ (\cref{sec:point-rasterizer}); any repair target that omits the $0.3$-pixel EWA dilation
is inconsistent with what the renderer will actually display.  For carrier $i$ with
contributor set $\mathcal{V}_i$ (from lineage) and observed 2D covariances $C_{iv}$, define
the whitened projection
\begin{equation}
  M_{iv} \;=\; C_{iv}^{-1/2}\,\big(J_{iv}\, \Sigma_i\, J_{iv}\T + 0.3\, I_2\big)\,
               C_{iv}^{-1/2},
  \label{eq:whitened}
\end{equation}
and minimize, over $\Sigma_i$ only (means, opacity, appearance, lineage, and count stay
bit-frozen), the symmetric generalized-log residual
\begin{equation}
  r_i \;=\; \sqrt{\;\frac{1}{|\mathcal{V}_i|}\sum_{v \in \mathcal{V}_i}
        \operatorname{mean}_k \log^2 \lambda_k(M_{iv})\;},
  \label{eq:cov-residual}
\end{equation}
with a Huber data term, a prior toward the Beam covariance, and explicit SPD/aspect bounds.
The $\log$-eigenvalue form is the load-bearing choice: it penalizes a projection twice as
large and half as large equally.  The previous repair minimized a one-sided whitened
Frobenius residual without the dilation term; that objective penalizes expansion more than
collapse, is shrink-biased, and was empirically harmful --- it worsened factorial-marginal
validation $\Jq$ by $51.6\%$.  The corrected objective of
\cref{eq:whitened,eq:cov-residual} reduced validation $\Jq$ by $63.2\%$ versus directly
optimizing raw Beam output, with 3/3 paired-root wins.\evidence{ara/logic/claims.md C30}
Because some fitted targets are sharper than the renderer's $0.3$-pixel floor, zero residual
is not always attainable; the repair aims at consistency, not perfection.

\paragraph{What is deliberately \emph{not} repaired.}
Opacity: a fitted 2D amplitude is a normalized-blend gauge quantity
(\cref{sec:field-model}), not an identifiable per-carrier 3D alpha under compositing ---
legacy opacity repair worsened $\Jq$ by $9.8\%$ and is retired.  Appearance: cross-view
amplitude-weighted color averaging has the same gauge problem and was immaterial ($0.12\%$,
below the $5\%$ necessity floor); SH0 colors are instead co-adapted by the compact phases.
Topology: even an optical-density-preserving clone (copying opacity changes coincident alpha
from $a$ to $1-(1-a)^2$; the preserving variant uses $a' = 1 - \sqrt{1-a}$) failed its stage
gate; the carrier path grows nothing.\evidence{ara/logic/claims.md C30;
docs/ADR-002-carrier-refinement.md}

\subsection{Two compact phases; why all parameter families train together}
\label{sec:phases}

Both phases minimize the sampled compact objective of \cref{eq:sampled-loss} at fixed
topology and SH degree 0.  Phase 1 trains all families jointly: the alpha/color gauge freedom
of the compositor means opacity and support cannot be identified independently --- scale,
orientation, opacity, and color must co-adapt.  Means-only optimization was strongly
rejected ($4.18\times$ the all-family phase-1 $\Jq$); a second all-family phase reduced
$\Jq$ to $0.717\times$ its stopping value with 3/3 wins; freezing \emph{only the means} in
phase 2 costs $1.47\%$ $\Jq$ and passes the frozen $2\%$ non-inferiority gate --- which is
exactly what makes the containment invariant of \cref{sec:containment} permanent.
Incremental SH3 improved $\Jq$ by only $2.3$--$3.0\%$, below the $5\%$ materiality floor, so
the carrier stays SH0.\evidence{ara/logic/claims.md C30}

\subsection{Projected-center containment}
\label{sec:containment}

Raw masks are forbidden inside the compact arms; their legal surrogate is the union of fitted
2D Gaussian support ellipses.  After phase 1, carrier $i$ is retained only if, for
\emph{every} fitting view $v$:
\begin{equation}
  z_{iv} > z_{\text{near}}
  \qquad\text{and}\qquad
  \min_j \big(\cam_v(\bm{\mu}_i) - \bm{m}_{vj}\big)\T C_{vj}^{-1}
         \big(\cam_v(\bm{\mu}_i) - \bm{m}_{vj}\big) \;\le\; 3^2,
  \label{eq:containment}
\end{equation}
i.e.\ the projected center must fall inside some positive-amplitude fitted component's
$3\sigma$ ellipse in every view that observes the scene.  The threshold is fixed and cannot
be weakened through configuration; phase 2 freezes means and rechecks
\cref{eq:containment}, so the invariant survives to the output.  On the development scene the
prune removes $5.3$--$5.4\%$ of carriers, retains $4{,}729$--$4{,}734$, ends with exactly
zero violations, and costs $2.69\%$ $\Jq$ against the unpruned
anchor.\evidence{ara/logic/claims.md C30, C31}

Honesty about what \cref{eq:containment} is: a \emph{visual-hull center invariant}.  It
removes carriers whose centers leave the fitted silhouette cone; it cannot detect a floater
\emph{inside} the multi-view hull, and it does not prove surface occupancy.  Whole-Gaussian
mask containment is impossible in principle (Gaussian support is infinite), and a
finite-footprint rule would erode legitimate silhouette-boundary primitives.  A
count-preserving soft variant is predeclared for the ablation grid: the barrier
\begin{equation}
  L_{\text{support}} = \sum_{i,v}
    \softplus\!\Big(\tfrac{\min_j q_{ivj} - 9}{\tau}\Big)^{2}
    + \lambda_z\, \softplus\!\Big(\tfrac{z_{\text{near}} - z_{iv}}{\tau_z}\Big)^{2},
  \label{eq:soft-support}
\end{equation}
optionally with deterministic samples on the projected $3\sigma$ ellipse to constrain the
footprint rather than the center.  \toshow{The tested nearest-ellipse hinge variant did not
reduce violations materially and is retired at that formulation; \cref{eq:soft-support} is a
design to freeze before use, not an evaluated method.  An independently measurable
surface-occupancy test for hull-interior floaters (depth reference, cross-view occupancy, or
geometry ground truth) remains an open obligation (O-12) --- ``no physical free-floaters''
must not be claimed from \cref{eq:containment} alone.}

\subsection{What the earlier RGB-backed schedules taught us}
\label{sec:carrier-history}

Two negative development results shaped this design and belong in the record.  (i) A short
fixed-budget carrier schedule (30/40/30/60 updates with clone waves, followed by a 40-update
recovery) \emph{preserved} its carriers ($94.6\%$ survival) yet finished $1.57$~dB below
immediate standard refinement; removing covariance repair \emph{improved} the endpoint by
$1.27$~dB even though the repair halved its own local residual ($-50.2\%$) --- local repair
objectives and downstream quality can dissociate when correspondences are
wrong.\evidence{ara/logic/claims.md C28}  (ii) The instrumented all-view maturation run
completed but violated its convergence-between-stages contract: three clone recoveries and
the higher-SH phase exhausted their 15k-update caps without reaching the plateau
criterion.\evidence{ara/logic/claims.md C29}  Both runs consumed RGB during optimization.
The compact-only sequence above replaces them: paired controls, predeclared budgets, no
topology growth, and a hard no-image boundary.  The dissociation lesson survives as a design
rule: \emph{every} local repair stage must justify itself by the matched no-repair ablation
on downstream risk (which the corrected repair now does, \cref{sec:cov-repair}), never by
its own residual.

\begin{figure}[t]
  \centering
  \figplaceholder{5.5cm}{%
    \textbf{Figure 9 --- carrier maturation stages (visual).}  One row per stage (Beam raw,
    after covariance repair, after phase 1, after prune, after phase 2): rendered front view
    + one side view, with carrier count and $\Jq$ annotated.  Purpose: make each accepted
    stage's contribution visible; the covariance-repair panel should visibly de-elongate
    beams.  Data: representative-root artifacts of the 2026-07-28 compact-only runs
    (\repopath{runs/}; re-render with the repository viewer).}
  \caption{\redcaption{Stage-by-stage carrier maturation renders.  To be produced from
    existing development artifacts.}}
  \label{fig:carrier-stages}
\end{figure}
